\scanpage{21}
% archive sheet number: 20
% L1
\deleted{Prenons une résolution injective $C^\cdot(F)$ de $F$,
% L2
alors $C^\cdot(F)_{S_\infty}$ est une résolution injective
% L3
de $\deleted{\unsure{C}}F_{S_\infty}$, et \deleted{\unsure{le|les}} $C^\cdot(F)_{S_i}=C^\cdot(F)_{S_\infty}\mid S_i$
% L4
une résolution injective de $F_{S_i}=F_{S_\infty}\mid S_i$.}

% L5
\deleted{On a}
% L6
\[
\deleted{\Gamma(S_\infty,C^\cdot(F)_{S_\infty})=\varprojlim\nolimits_i\Gamma(S_i,C^\cdot(F)_{S_i})}
\]
% L7
\deleted{le système projectif du $2^{\mathrm e}$ \unsure{mb} étant \uline{strict}.
% L8
\deleted{et} Par suite, \deleted{\unsure{utilisant}}}
% L9
\[
\deleted{H^n(S_\infty,F)\longrightarrow\varprojlim\nolimits_i H^n(S_i,F)}
\]
% L10
\deleted{et cet \unsure{hom|bien} est \uline{surjectif}\,; il est
% L11
bijectif si le système des $H^{n-1}(S_i,F)$
% L12
vérifie ML. Dans le cas général, il
% L13
semble qu'il faille écrire}
% L14
\[
\deleted{H^*(S_\infty,F)\Longleftarrow E^{pq}_2=\varprojlim\nolimits_i^{(p)} H^q(S_i,F)}
\]
% L15
\deleted{mais avec $\varprojlim{}^{(p)}=0$ si $p\ne 0,1$, donc des
% L16
suites exactes}
% L17
\[
\deleted{0\longrightarrow\varprojlim\nolimits_i^{(1)}H^{n-1}(S_i,F)\to H^n(S_\infty,F)\to\varprojlim\nolimits_i H^n(S_i,F)\to 0}
\]

% L18
\deleted{Dans le cas qui nous occupe, il vient :}

